\RequirePackage{plautopatch}
\documentclass[platex,dvipdfmx,a4paper]{jsarticle}

\usepackage{amsmath}
\usepackage{amsfonts}
\usepackage{amssymb}

\usepackage[dvipdfmx]{graphicx}
\usepackage{float}

\usepackage{tikz}
\usepackage{caption}

% 「%」は以降の内容を「改行コードも含めて」無視するコマンド
\usepackage[%
  dvipdfmx,% 欧文ではコメントアウトする
  setpagesize=false,%
  bookmarks=true,%
  bookmarksdepth=tocdepth,%
  bookmarksnumbered=true,%
  colorlinks=false,%
  pdftitle={},%
  pdfsubject={},%
  pdfauthor={},%
  pdfkeywords={}%
]{hyperref}
% PDFのしおり機能の日本語文字化けを防ぐ((u)pLaTeXのときのみかく)
\usepackage{pxjahyper}

\usepackage{framed}
\setlength{\OuterFrameSep}{0pt}

\usepackage{physics2}
\usephysicsmodule{ab}
\usephysicsmodule{ab.braket}


\newcommand{\hb}[2]{\frac{\partial #1}{\partial #2}}
\newcommand{\bbn}[2]{\frac{d #1}{d #2}}
\newcommand{\nhb}[3]{\frac{\partial^{#1} #2}{\partial #3^{#1}}}
\newcommand{\nbbn}[3]{\frac{d^{#1} #2}{d #3^{#1}}}
\newcommand{\vectxy}[2]{\begin{bmatrix}#1\\#2\end{bmatrix}}
\newcommand{\vectxyz}[3]{\begin{bmatrix}#1\\#2\\#3\end{bmatrix}}
\newcommand{\vect}[1]{\boldsymbol{\mathrm{#1}}}

\newcommand{\ssum}[3]{\sum_{#1=#2}^#3}

\newcommand{\veff}{V_{\mathrm{eff}}(r)}

\hypersetup{colorlinks=true}
\hypersetup{linkcolor=red}

\title{解析力学演習レポート2}
\author{25B00729　嶋田 悠音}

\begin{document}
\maketitle
\subsection*{2}

電気双極子放射の導出では，電荷分布の大きさ $d$ に対して変化の時間尺度 $\tau$ が十分長く，

\begin{align*}
\frac{d}{c\tau}\ll 1
\end{align*}

であるとして遅延時間の効果を近似した。ここではその仮定を外し，点電荷

\begin{align*}
\rho(\vect{r},t)
=
q\delta^3(\vect{r}-\vect{s}(t))
\end{align*}

に対する厳密な遅延ポテンシャルを求める。

与えられた式に空間積分を適用すると，

\begin{align*}
\phi(\vect{x},t)
&=
\frac{q}{4\pi\varepsilon_0}
\int dt''
\frac{\delta\!\left(
t-t''-\frac{R(t'')}{c}
\right)}
{R(t'')} ,
\end{align*}

ただし

\begin{align*}
R(t'')
=
|\vect{x}-\vect{s}(t'')|
\end{align*}

とおいた。

ここで

\begin{align*}
f(t'')
=
t-t''-\frac{R(t'')}{c}
\end{align*}

と定義すると，遅延時刻 $t_r$ は

\begin{align*}
f(t_r)=0
\end{align*}

すなわち

\begin{align*}
t-t_r
=
\frac{R(t_r)}{c}
\end{align*}

を満たす。

デルタ関数の公式

\begin{align*}
\delta(f(t''))
=
\frac{\delta(t''-t_r)}
{|f'(t_r)|}
\end{align*}

を用いるため，まず $f'(t'')$ を求める。

\begin{align*}
f'(t'')
=
-1-\frac{1}{c}\frac{dR}{dt''}.
\end{align*}

一方，

\begin{align*}
R
=
|\vect{x}-\vect{s}(t'')|
\end{align*}

より，

\begin{align*}
\frac{dR}{dt''}
&=
\frac{\vect{x}-\vect{s}}
{|\vect{x}-\vect{s}|}
\cdot
\frac{d(\vect{x}-\vect{s})}{dt''} \\
&=
-\hat{\vect{R}}\cdot \vect{v},
\end{align*}

ただし

\begin{align*}
\hat{\vect{R}}
=
\frac{\vect{x}-\vect{s}(t'')}
     {|\vect{x}-\vect{s}(t'')|}
,\qquad
\vect{v}
=
\dot{\vect{s}}(t'')
\end{align*}

である。

したがって，

\begin{align*}
f'(t_r)
=
-1+\frac{\hat{\vect{R}}\cdot\vect{v}}{c}.
\end{align*}

よって

\begin{align*}
|f'(t_r)|
=
1-\frac{\hat{\vect{R}}\cdot\vect{v}}{c}
\end{align*}

となり，

\begin{align*}
\phi(\vect{x},t)
&=
\frac{q}{4\pi\varepsilon_0}
\frac{1}
{R(t_r)\left(
1-\dfrac{\hat{\vect{R}}\cdot\vect{v}}{c}
\right)}.
\end{align*}

これが Liénard-Wiechert スカラーポテンシャルである。

同様にベクトルポテンシャルは

\begin{align*}
\vect{A}(\vect{x},t)
=
\frac{\mu_0 q}{4\pi}
\frac{\vect{v}(t_r)}
{R(t_r)\left(
1-\dfrac{\hat{\vect{R}}\cdot\vect{v}}{c}
\right)}
\end{align*}

となる。

この結果は，多重極展開が本質的に低速・長波長近似であることを示している。電気双極子放射の導出では，電荷分布全体に対する遅延時間

\begin{align*}
\frac{|\vect{r}'|}{c}
\end{align*}

を小さい量として展開した。しかし電荷の運動速度が大きくなる場合や，系の大きさが放射波長と同程度になる場合には，空間の各点から到達する信号の遅延を無視できない。

Liénard-Wiechert ポテンシャルは，各電荷の寄与をそれぞれの遅延時刻 $t_r$ で評価した厳密解であり，多重極展開はこの厳密解を

\begin{align*}
\frac{d}{\lambda}\ll 1
\end{align*}

のもとで展開した近似に対応する。したがって高周波放射や相対論的運動では，単純な双極子近似ではなく Liénard-Wiechert ポテンシャルから電磁場を求める必要がある。

以上より，時間変化が十分遅いという仮定を外すと，ポテンシャルは各電荷の遅延時刻を用いて評価される Liénard-Wiechert ポテンシャルで記述され，多重極展開はその長波長近似として理解されることが示された。

\end{document}